\clearpage
\begin{explanationbox}
\textbf{\textcolor{CExplanation}{一、为什么会出现这条规律？}}
\begin{center}\begin{tikzpicture}[font=\small,>=stealth,line join=round]
\draw[gray!40,rounded corners=3pt] (0,0) rectangle (14.2,8.8);\node at (7.1,8.200000000000001) {奇配轴、偶配点：隔2变号，隔4回到原高度};
\node at (3.1500000000000004,7.3) {奇函数配轴 $x=1$};\begin{scope}[shift={(2.061111111111111,4.345714285714285)},x=0.5444444444444445cm,y=1.0857142857142856cm]\draw[->,gray] (-2.5,0)--(6.5,0) node[right]{$x$};\draw[->,gray] (0,-1.7)--(0,1.8) node[above]{$y$};\draw[blue,thick,solid,domain=-2:-1,samples=180] plot(\x,{-1*\x+(-2)});\draw[blue,thick,solid,domain=-1:0,samples=180] plot(\x,{1*\x+(0)});\draw[blue,thick,solid,domain=0:1,samples=180] plot(\x,{1*\x+(0)});\draw[blue,thick,solid,domain=1:2,samples=180] plot(\x,{-1*\x+(2)});\draw[blue,thick,solid,domain=2:3,samples=180] plot(\x,{-1*\x+(2)});\draw[blue,thick,solid,domain=3:4,samples=180] plot(\x,{1*\x+(-4)});\draw[blue,thick,solid,domain=4:5,samples=180] plot(\x,{1*\x+(-4)});\draw[blue,thick,solid,domain=5:6,samples=180] plot(\x,{-1*\x+(6)});\fill[orange!85!black] (1,1) circle(2pt);\fill[orange!85!black] (3,-1) circle(2pt);\fill[orange!85!black] (5,1) circle(2pt);\draw[violet,dashed] (1,-1.4)--(1,1.4);\end{scope}\node[align=center,text width=5.8cm] at (3.1500000000000004,1.7) {输入 $1,3,5$ 对应 $1,-1,1$\\一个正周期为4};\node at (9.95,7.3) {偶函数配中心 $(1,0)$};\begin{scope}[shift={(8.86111111111111,4.345714285714285)},x=0.5444444444444445cm,y=1.0857142857142856cm]\draw[->,gray] (-2.5,0)--(6.5,0) node[right]{$x$};\draw[->,gray] (0,-1.7)--(0,1.8) node[above]{$y$};\draw[teal,thick,solid,domain=-2:-1,samples=180] plot(\x,{1*\x+(1)});\draw[teal,thick,solid,domain=-1:0,samples=180] plot(\x,{1*\x+(1)});\draw[teal,thick,solid,domain=0:1,samples=180] plot(\x,{-1*\x+(1)});\draw[teal,thick,solid,domain=1:2,samples=180] plot(\x,{-1*\x+(1)});\draw[teal,thick,solid,domain=2:3,samples=180] plot(\x,{1*\x+(-3)});\draw[teal,thick,solid,domain=3:4,samples=180] plot(\x,{1*\x+(-3)});\draw[teal,thick,solid,domain=4:5,samples=180] plot(\x,{-1*\x+(5)});\draw[teal,thick,solid,domain=5:6,samples=180] plot(\x,{-1*\x+(5)});\fill[orange!85!black] (0,1) circle(2pt);\fill[orange!85!black] (2,-1) circle(2pt);\fill[orange!85!black] (4,1) circle(2pt);\fill[orange!85!black] (1,0) circle(2pt);\end{scope}\node[align=center,text width=5.8cm] at (9.95,1.7) {输入 $0,2,4$ 对应 $1,-1,1$\\一个正周期为4};
\node[text width=12.6cm,align=center] at (7.1,.65) {两图同尺度，均为真实分段直线函数；高低相反时还不能叫周期。};
\end{tikzpicture}\end{center}



以前学过：奇函数图像关于\textbf{原点}对称，偶函数图像关于\textbf{$y$ 轴}对称。现在题目又额外送给我们一条对称轴 $x=a$，或者一个对称中心 $(a,0)$。\textbf{同一张图，既能按第一种办法翻过去，又能按第二种办法翻回来}，于是便可以把远处的函数值和原来的值联系起来。

可以把它想成一排有正负号的数字：某个位置是 $3$，移动一段距离可能变成 $-3$，再移动同样距离又变成 $3$。\textbf{正、负、正、负……，两次之后恢复原样。}这解释的是奇配轴、偶配点为何给出 $4|a|$；偶配轴、奇配点则移动 $2|a|$ 后高度直接恢复。

\medskip

\end{explanationbox}
\clearpage
\begin{explanationbox}
\textbf{\textcolor{CExplanation}{二、先不列方程：用初中几何的“翻折”来理解}}
\begin{center}\begin{tikzpicture}[font=\small,>=stealth,line join=round]
\draw[gray!40,rounded corners=3pt] (0,0) rectangle (14.2,8.3);\node at (7.1,7.700000000000001) {证明一：奇配轴，追踪同一点的四个对应位置};
\begin{scope}[shift={(2.7666666666666666,3.8842857142857143)},x=0.6666666666666666cm,y=1.3142857142857143cm]\draw[->,gray] (-2.5,0)--(6.5,0) node[right]{$x$};\draw[->,gray] (0,-1.7)--(0,1.8) node[above]{$y$};\draw[blue,thick,solid,domain=-2:-1,samples=180] plot(\x,{-1*\x+(-2)});\draw[blue,thick,solid,domain=-1:0,samples=180] plot(\x,{1*\x+(0)});\draw[blue,thick,solid,domain=0:1,samples=180] plot(\x,{1*\x+(0)});\draw[blue,thick,solid,domain=1:2,samples=180] plot(\x,{-1*\x+(2)});\draw[blue,thick,solid,domain=2:3,samples=180] plot(\x,{-1*\x+(2)});\draw[blue,thick,solid,domain=3:4,samples=180] plot(\x,{1*\x+(-4)});\draw[blue,thick,solid,domain=4:5,samples=180] plot(\x,{1*\x+(-4)});\draw[blue,thick,solid,domain=5:6,samples=180] plot(\x,{-1*\x+(6)});\fill[orange!85!black] (0.5,0.5) circle(2.2pt);\fill[orange!85!black] (-0.5,-0.5) circle(2.2pt);\fill[orange!85!black] (2.5,-0.5) circle(2.2pt);\fill[orange!85!black] (4.5,0.5) circle(2.2pt);\foreach \X/\Y/\L in {.5/.5/P,-.5/-0.5/Q,2.5/-.5/R,4.5/.5/S}{\draw[gray,dashed] (\X,\Y)--(\X,-1.15);\node[below=4pt] at (\X,-1.45){$\L$};}\end{scope}\node[anchor=west,align=left,text width=5.0cm] at (8.4,4.0) {$P:(\frac12,\frac12)$\\[9pt]$Q:(-\frac12,-\frac12)$\\[9pt]$R:(\frac52,-\frac12)$\\[9pt]$S:(\frac92,\frac12)$\\[9pt]$P,S$ 横向相距4};
\node[text width=12.6cm,align=center] at (7.1,.65) {点的位置和高度都由折线计算；第一组对称得到相反高度，再做一组还原。};
\end{tikzpicture}\end{center}



\textbf{情况 A：奇函数 $+$ 关于 $x=a$ 对称。}从图像上任取一点 $P(x,y)$：
\begin{enumerate}[leftmargin=*,itemsep=0.25em]
 \item 奇函数说明图像绕原点转半圈仍是原图，所以有点 $(-x,-y)$；
 \item 再把这个点沿直线 $x=a$ 翻折，左右位置距 $a$ 相同，得到 $(2a+x,-y)$；
 \item 现在横坐标从 $x$ 变成 $x+2a$，但纵坐标从 $y$ 变成 $-y$；
 \item 再做同样两步，纵坐标从 $-y$ 变回 $y$，横坐标从 $x+2a$ 变成 $x+4a$。
\end{enumerate}
即
\[
(x,y)\longrightarrow(x+2a,-y)\longrightarrow(x+4a,y).
\]

\end{explanationbox}\clearpage\begin{explanationbox}
\begin{center}\begin{tikzpicture}[font=\small,>=stealth,line join=round]
\draw[gray!40,rounded corners=3pt] (0,0) rectangle (14.2,8.3);\node at (7.1,7.700000000000001) {证明一：偶配点，追踪同一点的四个对应位置};
\begin{scope}[shift={(2.7666666666666666,3.8842857142857143)},x=0.6666666666666666cm,y=1.3142857142857143cm]\draw[->,gray] (-2.5,0)--(6.5,0) node[right]{$x$};\draw[->,gray] (0,-1.7)--(0,1.8) node[above]{$y$};\draw[teal,thick,solid,domain=-2:-1,samples=180] plot(\x,{1*\x+(1)});\draw[teal,thick,solid,domain=-1:0,samples=180] plot(\x,{1*\x+(1)});\draw[teal,thick,solid,domain=0:1,samples=180] plot(\x,{-1*\x+(1)});\draw[teal,thick,solid,domain=1:2,samples=180] plot(\x,{-1*\x+(1)});\draw[teal,thick,solid,domain=2:3,samples=180] plot(\x,{1*\x+(-3)});\draw[teal,thick,solid,domain=3:4,samples=180] plot(\x,{1*\x+(-3)});\draw[teal,thick,solid,domain=4:5,samples=180] plot(\x,{-1*\x+(5)});\draw[teal,thick,solid,domain=5:6,samples=180] plot(\x,{-1*\x+(5)});\fill[orange!85!black] (0.5,0.5) circle(2.2pt);\fill[orange!85!black] (-0.5,0.5) circle(2.2pt);\fill[orange!85!black] (2.5,-0.5) circle(2.2pt);\fill[orange!85!black] (4.5,0.5) circle(2.2pt);\foreach \X/\Y/\L in {.5/.5/P,-.5/0.5/Q,2.5/-.5/R,4.5/.5/S}{\draw[gray,dashed] (\X,\Y)--(\X,-1.15);\node[below=4pt] at (\X,-1.45){$\L$};}\end{scope}\node[anchor=west,align=left,text width=5.0cm] at (8.4,4.0) {$P:(\frac12,\frac12)$\\[9pt]$Q:(-\frac12,\frac12)$\\[9pt]$R:(\frac52,-\frac12)$\\[9pt]$S:(\frac92,\frac12)$\\[9pt]$P,S$ 横向相距4};
\node[text width=12.6cm,align=center] at (7.1,.65) {点的位置和高度都由折线计算；第一组对称得到相反高度，再做一组还原。};
\end{tikzpicture}\end{center}

\textbf{情况 B：偶函数 $+$ 关于点 $(a,0)$ 对称。}还是从 $P(x,y)$ 出发：
\begin{enumerate}[leftmargin=*,itemsep=0.25em]
 \item 偶函数先把点沿 $y$ 轴翻折，得到 $(-x,y)$，高度不变；
 \item 再以点 $(a,0)$ 为中心旋转半圈，得到 $(2a+x,-y)$；
 \item 一次走到 $x+2a$ 时高度变号；再做一次，走到 $x+4a$ 时高度恢复。
\end{enumerate}
结果同样是
\[
(x,y)\longrightarrow(x+2a,-y)\longrightarrow(x+4a,y).
\]

\textbf{一句话记忆：}有一“翻”改变上下，又有一“翻”改变左右；两种对称接起来，先反号，后恢复。

\medskip

\end{explanationbox}
\clearpage
\begin{explanationbox}
\textbf{\textcolor{CExplanation}{三、最容易看懂的例子：不学三角函数也可以}}
\begin{center}\begin{tikzpicture}[font=\small,>=stealth,line join=round]
\draw[gray!40,rounded corners=3pt] (0,0) rectangle (14.2,8.3);\node at (7.1,7.700000000000001) {例1：从第一段直线出发，按对称逐段拼出图像};
\begin{scope}[shift={(2.7666666666666666,3.8842857142857143)},x=0.6666666666666666cm,y=1.3142857142857143cm]\draw[->,gray] (-2.5,0)--(6.5,0) node[right]{$x$};\draw[->,gray] (0,-1.7)--(0,1.8) node[above]{$y$};\draw[blue,thick,solid,domain=-2:-1,samples=180] plot(\x,{-1*\x+(-2)});\draw[blue,thick,solid,domain=-1:0,samples=180] plot(\x,{1*\x+(0)});\draw[blue,thick,solid,domain=0:1,samples=180] plot(\x,{1*\x+(0)});\draw[blue,thick,solid,domain=1:2,samples=180] plot(\x,{-1*\x+(2)});\draw[blue,thick,solid,domain=2:3,samples=180] plot(\x,{-1*\x+(2)});\draw[blue,thick,solid,domain=3:4,samples=180] plot(\x,{1*\x+(-4)});\draw[blue,thick,solid,domain=4:5,samples=180] plot(\x,{1*\x+(-4)});\draw[blue,thick,solid,domain=5:6,samples=180] plot(\x,{-1*\x+(6)});\fill[orange!85!black] (0.5,0.5) circle(2.2pt);\fill[orange!85!black] (2.5,-0.5) circle(2.2pt);\fill[orange!85!black] (0,0) circle(2.2pt);\fill[orange!85!black] (1,1) circle(2.2pt);\fill[orange!85!black] (2,0) circle(2.2pt);\fill[orange!85!black] (3,-1) circle(2.2pt);\fill[orange!85!black] (4,0) circle(2.2pt);\foreach \v in {-2,-1,1,2,3,4,5,6}{\node[below=4pt] at (\v,-1.35){$\v$};}\end{scope}\node[anchor=west,align=left,text width=5.0cm] at (8.4,4.0) {$a=1$，奇函数\\[9pt]轴为 $x=1$\\[9pt]$f(\frac12)=\frac12$\\[9pt]$f(\frac52)=-\frac12$\\[9pt]$f(2026.5)=-\frac12$};
\node[text width=12.6cm,align=center] at (7.1,.65) {每段长度1；0到1先上升，1到2下降，再用反号关系补齐2到4。};
\end{tikzpicture}\end{center}



取 $a=1$，先在 $0\le x\le1$ 上画一段很普通的直线 $y=x$，从 $(0,0)$ 走到 $(1,1)$。接着按“关于 $x=1$ 对称”把它翻到右侧，再按“奇函数关于原点对称”把左右接起来；继续重复，得到一条由直线段组成的上下起伏的图像。

\begin{center}
\begin{tabular}{c|rrrrrrrrr}
$x$ & $-2$ & $-1$ & $0$ & $1$ & $2$ & $3$ & $4$ & $5$ & $6$\\ \hline
$f(x)$ & $0$ & $-1$ & $0$ & $1$ & $0$ & $-1$ & $0$ & $1$ & $0$
\end{tabular}
\end{center}
你会发现：从 $x=1$ 的高点到 $x=3$ 的低点，距离是 $2$；从高点 $x=1$ 到下一个高点 $x=5$，距离是 $4$。这才是“\textbf{走两格变负，走四格还原}”的实际画面。

\end{explanationbox}\clearpage\begin{explanationbox}另一种搭配也能只用直线来画：让 $f(0)=1,f(1)=0,f(2)=-1,f(3)=0,f(4)=1$，相邻点用线段连起来，再按每隔 $4$ 重复。这条图像关于 $y$ 轴对称，也关于点 $(1,0)$ 对称。下图直接画出第二种搭配。
\begin{center}\begin{tikzpicture}[font=\small,>=stealth,line join=round]
\draw[gray!40,rounded corners=3pt] (0,0) rectangle (14.2,8.3);\node at (7.1,7.700000000000001) {例2：偶折线在0和4取1，在2处取负1};
\begin{scope}[shift={(2.7666666666666666,3.8842857142857143)},x=0.6666666666666666cm,y=1.3142857142857143cm]\draw[->,gray] (-2.5,0)--(6.5,0) node[right]{$x$};\draw[->,gray] (0,-1.7)--(0,1.8) node[above]{$y$};\draw[teal,thick,solid,domain=-2:-1,samples=180] plot(\x,{1*\x+(1)});\draw[teal,thick,solid,domain=-1:0,samples=180] plot(\x,{1*\x+(1)});\draw[teal,thick,solid,domain=0:1,samples=180] plot(\x,{-1*\x+(1)});\draw[teal,thick,solid,domain=1:2,samples=180] plot(\x,{-1*\x+(1)});\draw[teal,thick,solid,domain=2:3,samples=180] plot(\x,{1*\x+(-3)});\draw[teal,thick,solid,domain=3:4,samples=180] plot(\x,{1*\x+(-3)});\draw[teal,thick,solid,domain=4:5,samples=180] plot(\x,{-1*\x+(5)});\draw[teal,thick,solid,domain=5:6,samples=180] plot(\x,{-1*\x+(5)});\fill[orange!85!black] (0,1) circle(2.2pt);\fill[orange!85!black] (1,0) circle(2.2pt);\fill[orange!85!black] (2,-1) circle(2.2pt);\fill[orange!85!black] (3,0) circle(2.2pt);\fill[orange!85!black] (4,1) circle(2.2pt);\foreach \v in {-2,-1,1,2,3,4,5,6}{\node[below=4pt] at (\v,-1.35){$\v$};}\end{scope}\node[anchor=west,align=left,text width=5.0cm] at (8.4,4.0) {$g$ 为偶函数\\[9pt]中心 $(1,0)$\\[9pt]$[0,2]$ 上 $g=1-x$\\[9pt]$g(2)=-1$\\[9pt]$g(2026)=g(2)=-1$};
\node[text width=12.6cm,align=center] at (7.1,.65) {关于y轴等高，关于点(1,0)相反；两种对称一起使得隔2变号。};
\end{tikzpicture}\end{center}



\medskip

\end{explanationbox}
\clearpage
\begin{explanationbox}
\textbf{\textcolor{CExplanation}{四、为什么不是移动 $2|a|$ 就算周期？}}
\begin{center}\begin{tikzpicture}[font=\small,>=stealth,line join=round]
\draw[gray!40,rounded corners=3pt] (0,0) rectangle (14.2,8.3);\node at (7.1,7.700000000000001) {易错二：5先变负5，再回到5；隔2尚未还原};
\begin{scope}[shift={(1.836842105263158,3.8706896551724133)},x=1.0526315789473684cm,y=0.3172413793103448cm]\draw[->,gray] (-0.7,0)--(5,0) node[right]{$x$};\draw[->,gray] (0,-7)--(0,7.5) node[above]{$y$};\draw[teal,thick,solid,domain=-0.3:4.6,samples=180] plot(\x,{5*cos(deg(pi*\x/2))});\fill[orange!85!black] (0,5) circle(2.2pt);\fill[orange!85!black] (2,-5) circle(2.2pt);\fill[orange!85!black] (4,5) circle(2.2pt);\end{scope}\node[anchor=west,align=left,text width=5.0cm] at (8.4,4.0) {示例 $f=5\cos\frac{\pi x}2$\\[9pt]偶函数，中心 $(1,0)$\\[9pt]$f(0)=5$\\[9pt]$f(2)=-5\ne5$\\[9pt]$f(4)=5$};
\node[text width=12.6cm,align=center] at (7.1,.65) {本图直接使用误区中的数值5；反号不等于相等，非零函数不能把2当周期。};
\end{tikzpicture}\end{center}



以下专指奇配轴、偶配点两种搭配。“周期”要求移动后\textbf{每个位置的函数值都不变}，但这两种搭配先只能保证
\[f(x+2|a|)=-f(x).\]
例如原来是 $1$，移动后是 $-1$，并不相同；所以只要函数不是处处等于零，$2|a|$ 就\textbf{不是}周期。可这\textbf{不能}推出 $4|a|$ 一定最小：它们之间或更小的距离，也可能早已使整张图重复。求“最小正周期”还需额外信息。

\medskip

\end{explanationbox}
\clearpage
\begin{explanationbox}
\textbf{\textcolor{CExplanation}{五、和以前学过的知识连起来：四种对称组合}}
\begin{center}\begin{tikzpicture}[font=\small,>=stealth,line join=round]
\draw[gray!40,rounded corners=3pt] (0,0) rectangle (14.2,8.8);\node at (7.1,8.200000000000001) {同类对称只走2就还原，与奇配轴、偶配点不同};
\node at (3.1500000000000004,7.3) {奇配点：$f=\sin\pi x$};\begin{scope}[shift={(1.4795454545454545,4.338709677419355)},x=1.1136363636363635cm,y=1.225806451612903cm]\draw[->,gray] (-0.7,0)--(3.7,0) node[right]{$x$};\draw[->,gray] (0,-1.5)--(0,1.6) node[above]{$y$};\draw[blue,thick,solid,domain=-0.4:3.4,samples=180] plot(\x,{sin(deg(pi*\x))});\fill[orange!85!black] (0,0) circle(2pt);\fill[orange!85!black] (1,0) circle(2pt);\end{scope}\node[align=center,text width=5.8cm] at (3.1500000000000004,1.7) {中心 $(0,0),(1,0)$\\一个周期为2};\node at (9.95,7.3) {偶配轴：$g=\cos\pi x$};\begin{scope}[shift={(8.279545454545454,4.338709677419355)},x=1.1136363636363635cm,y=1.225806451612903cm]\draw[->,gray] (-0.7,0)--(3.7,0) node[right]{$x$};\draw[->,gray] (0,-1.5)--(0,1.6) node[above]{$y$};\draw[teal,thick,solid,domain=-0.4:3.4,samples=180] plot(\x,{cos(deg(pi*\x))});\foreach \a in {0,1}{\draw[violet,dashed] (\a,-1.3)--(\a,1.3);}\end{scope}\node[align=center,text width=5.8cm] at (9.95,1.7) {轴 $x=0,1$\\一个周期也为2};
\node[text width=12.6cm,align=center] at (7.1,.65) {两图同尺度：同类组合两次都不变号或两次变号，移2就保持原高度。};
\end{tikzpicture}\end{center}



\begin{center}
\begin{tabular}{p{0.27\linewidth}|p{0.30\linewidth}|p{0.29\linewidth}}
原来的性质 & 新增加的对称 & 可以直接推出\\ \hline
奇函数（原点对称） & 关于 $x=a$ 对称 & $f(x+2a)=-f(x)$，周期 $4|a|$\\
偶函数（$y$ 轴对称） & 关于 $(a,0)$ 对称 & $f(x+2a)=-f(x)$，周期 $4|a|$\\
奇函数（原点对称） & 关于 $(a,0)$ 对称 & $f(x+2a)=f(x)$，周期 $2|a|$\\
偶函数（$y$ 轴对称） & 关于 $x=a$ 对称 & $f(x+2a)=f(x)$，周期 $2|a|$
\end{tabular}
\end{center}
\textbf{眼前一亮的联系：}两种“同类”对称（两个点、两条轴），走 $2|a|$ 就重复；两种“一点一轴”对称，走 $2|a|$ 先变号，再走同样的距离才重复。这也能帮助区别相关结论 F026。

\medskip

\end{explanationbox}
\clearpage
\begin{explanationbox}
\textbf{\textcolor{CExplanation}{六、顺手多得两个函数值}}
\begin{center}\begin{tikzpicture}[font=\small,>=stealth,line join=round]
\draw[gray!40,rounded corners=3pt] (0,0) rectangle (14.2,8.8);\node at (7.1,8.200000000000001) {先看哪一种对称强制取0，再沿反号关系继续找};
\node at (3.1500000000000004,7.3) {奇配轴：$f=\sin\frac{\pi x}2$};\begin{scope}[shift={(1.293939393939394,4.338709677419355)},x=0.7424242424242425cm,y=1.225806451612903cm]\draw[->,gray] (-0.8,0)--(5.8,0) node[right]{$x$};\draw[->,gray] (0,-1.5)--(0,1.6) node[above]{$y$};\draw[blue,thick,solid,domain=-0.5:5.5,samples=180] plot(\x,{sin(deg(pi*\x/2))});\fill[orange!85!black] (0,0) circle(2pt);\fill[orange!85!black] (2,0) circle(2pt);\fill[orange!85!black] (4,0) circle(2pt);\end{scope}\node[align=center,text width=5.8cm] at (3.1500000000000004,1.7) {奇函数先给 $f(0)=0$\\再得到 $f(2)=f(4)=0$};\node at (9.95,7.3) {偶配点：$g=\cos\frac{\pi x}2$};\begin{scope}[shift={(8.093939393939394,4.338709677419355)},x=0.7424242424242425cm,y=1.225806451612903cm]\draw[->,gray] (-0.8,0)--(5.8,0) node[right]{$x$};\draw[->,gray] (0,-1.5)--(0,1.6) node[above]{$y$};\draw[teal,thick,solid,domain=-0.5:5.5,samples=180] plot(\x,{cos(deg(pi*\x/2))});\fill[orange!85!black] (1,0) circle(2pt);\fill[orange!85!black] (3,0) circle(2pt);\fill[orange!85!black] (5,0) circle(2pt);\end{scope}\node[align=center,text width=5.8cm] at (9.95,1.7) {中心先给 $g(1)=0$\\再得到 $g(3)=g(5)=0$};
\node[text width=12.6cm,align=center] at (7.1,.65) {两图同尺度；奇配轴不能直接断言轴上的函数值为0，例如f(1)=1。};
\end{tikzpicture}\end{center}



若是奇函数又关于 $x=a$ 对称，那么 $f(0)=0$，并由 $f(2a)=-f(0)=0$ 得到 $f(2a)=0$；进一步每个 $x=2na$（$n$ 为整数）处都是 $0$。

若是偶函数又关于 $(a,0)$ 对称，把对称点本身代入得 $f(a)=-f(a)$，所以 $f(a)=0$；再由每隔 $2a$ 变号，可知 $f((2n+1)a)=0$（$n$ 为整数）。\textbf{不用函数式，就能锁定整排零点。}
\end{explanationbox}
\clearpage
\begin{explanationbox}
\textbf{四图逐项看｜偶函数 + 轴 $x=1$}
\begin{center}\begin{tikzpicture}[font=\small,>=stealth,line join=round]
\draw[gray!40,rounded corners=3pt] (0,0) rectangle (14.2,8.3);
\node[text width=12.6cm,align=center] at (7.1,7.65) {偶函数 + 轴 $x=1$};

\begin{scope}[shift={(3.2,4.3)},x=1.35cm,y=1.15cm]
\draw[->,gray] (-1.7,0)--(5.25,0) node[right]{$x$};
\draw[->,gray] (0,-1.45)--(0,1.65) node[above]{$y$};
\draw[blue,thick,domain=-1.5:5,samples=240] plot(\x,{cos(deg(pi*\x))});
\draw[violet,dashed] (1,-1.4)--(1,1.45) node[above]{$x=1$};
\fill[orange!85!black] (-0.25,{cos(deg(pi*-0.25))}) circle(2.2pt);
\fill[orange!85!black] (0.25,{cos(deg(pi*0.25))}) circle(2.2pt);
\fill[orange!85!black] (2.25,{cos(deg(pi*2.25))}) circle(2.2pt);
\fill[orange!85!black] (4.25,{cos(deg(pi*4.25))}) circle(2.2pt);
\foreach \v in {-1,1,2,3,4,5}{\node[below=4pt] at (\v,-1.5){$\v$};}\end{scope}
\node at (7.1,6.85) {$f(x)=\cos(\pi x),\quad a=1$};
\node[align=center] at (7.1,1.3) {$f(x+2)=f(x)$；$2$ 是一个正周期。};
\node[text width=12.6cm,align=center] at (7.1,.65) {橙点的横坐标为 $-\frac14,\frac14,\frac94,\frac{17}4$；曲线仅是满足全部条件的一个例子。};
\end{tikzpicture}\end{center}

偶函数本来就关于 $y$ 轴对称。点 $(x,y)$ 先沿 $y$ 轴翻到 $(-x,y)$，再沿 $x=a$ 翻到 $(2a+x,y)$。两次都只左右翻，高度一直是 $y$，所以移动 $2|a|$ 已经重复。
\end{explanationbox}
\clearpage
\begin{explanationbox}
\textbf{四图逐项看｜偶函数 + 中心 $(1,0)$}
\begin{center}\begin{tikzpicture}[font=\small,>=stealth,line join=round]
\draw[gray!40,rounded corners=3pt] (0,0) rectangle (14.2,8.3);
\node[text width=12.6cm,align=center] at (7.1,7.65) {偶函数 + 中心 $(1,0)$};

\begin{scope}[shift={(3.2,4.3)},x=1.35cm,y=1.15cm]
\draw[->,gray] (-1.7,0)--(5.25,0) node[right]{$x$};
\draw[->,gray] (0,-1.45)--(0,1.65) node[above]{$y$};
\draw[blue,thick,domain=-1.5:5,samples=240] plot(\x,{cos(deg(pi*\x/2))});
\fill[violet] (1,0) circle(2.8pt);\node[above right=4pt,violet] at (1,0){$(1,0)$};
\fill[orange!85!black] (-0.25,{cos(deg(pi*-0.25/2))}) circle(2.2pt);
\fill[orange!85!black] (0.25,{cos(deg(pi*0.25/2))}) circle(2.2pt);
\fill[orange!85!black] (2.25,{cos(deg(pi*2.25/2))}) circle(2.2pt);
\fill[orange!85!black] (4.25,{cos(deg(pi*4.25/2))}) circle(2.2pt);
\foreach \v in {-1,1,2,3,4,5}{\node[below=4pt] at (\v,-1.5){$\v$};}\end{scope}
\node at (7.1,6.85) {$f(x)=\cos(\pi x/2),\quad a=1$};
\node[align=center] at (7.1,1.3) {$f(x+2)=-f(x)$；$4$ 是一个正周期。};
\node[text width=12.6cm,align=center] at (7.1,.65) {橙点的横坐标为 $-\frac14,\frac14,\frac94,\frac{17}4$；曲线仅是满足全部条件的一个例子。};
\end{tikzpicture}\end{center}

偶函数先沿 $y$ 轴翻到 $(-x,y)$，新增的中心 $(a,0)$ 再把它变成 $(2a+x,-y)$。两种对称合起来只改变一次高度的正负；再做一轮才得到 $(x+4a,y)$。
\end{explanationbox}
\clearpage
\begin{explanationbox}
\textbf{四图逐项看｜奇函数 + 轴 $x=1$}
\begin{center}\begin{tikzpicture}[font=\small,>=stealth,line join=round]
\draw[gray!40,rounded corners=3pt] (0,0) rectangle (14.2,8.3);
\node[text width=12.6cm,align=center] at (7.1,7.65) {奇函数 + 轴 $x=1$};

\begin{scope}[shift={(3.2,4.3)},x=1.35cm,y=1.15cm]
\draw[->,gray] (-1.7,0)--(5.25,0) node[right]{$x$};
\draw[->,gray] (0,-1.45)--(0,1.65) node[above]{$y$};
\draw[blue,thick,domain=-1.5:5,samples=240] plot(\x,{sin(deg(pi*\x/2))});
\draw[violet,dashed] (1,-1.4)--(1,1.45) node[above]{$x=1$};
\fill[orange!85!black] (-0.25,{sin(deg(pi*-0.25/2))}) circle(2.2pt);
\fill[orange!85!black] (0.25,{sin(deg(pi*0.25/2))}) circle(2.2pt);
\fill[orange!85!black] (2.25,{sin(deg(pi*2.25/2))}) circle(2.2pt);
\fill[orange!85!black] (4.25,{sin(deg(pi*4.25/2))}) circle(2.2pt);
\foreach \v in {-1,1,2,3,4,5}{\node[below=4pt] at (\v,-1.5){$\v$};}\end{scope}
\node at (7.1,6.85) {$f(x)=\sin(\pi x/2),\quad a=1$};
\node[align=center] at (7.1,1.3) {$f(x+2)=-f(x)$；$4$ 是一个正周期。};
\node[text width=12.6cm,align=center] at (7.1,.65) {橙点的横坐标为 $-\frac14,\frac14,\frac94,\frac{17}4$；曲线仅是满足全部条件的一个例子。};
\end{tikzpicture}\end{center}

奇函数先绕原点转到 $(-x,-y)$，新增的轴 $x=a$ 再把它变成 $(2a+x,-y)$。高度只变号一次；再做一轮才得到 $(x+4a,y)$。
\end{explanationbox}
\clearpage
\begin{explanationbox}
\textbf{四图逐项看｜奇函数 + 中心 $(1,0)$}
\begin{center}\begin{tikzpicture}[font=\small,>=stealth,line join=round]
\draw[gray!40,rounded corners=3pt] (0,0) rectangle (14.2,8.3);
\node[text width=12.6cm,align=center] at (7.1,7.65) {奇函数 + 中心 $(1,0)$};

\begin{scope}[shift={(3.2,4.3)},x=1.35cm,y=1.15cm]
\draw[->,gray] (-1.7,0)--(5.25,0) node[right]{$x$};
\draw[->,gray] (0,-1.45)--(0,1.65) node[above]{$y$};
\draw[blue,thick,domain=-1.5:5,samples=240] plot(\x,{sin(deg(pi*\x))});
\fill[violet] (1,0) circle(2.8pt);\node[above right=4pt,violet] at (1,0){$(1,0)$};
\fill[orange!85!black] (-0.25,{sin(deg(pi*-0.25))}) circle(2.2pt);
\fill[orange!85!black] (0.25,{sin(deg(pi*0.25))}) circle(2.2pt);
\fill[orange!85!black] (2.25,{sin(deg(pi*2.25))}) circle(2.2pt);
\fill[orange!85!black] (4.25,{sin(deg(pi*4.25))}) circle(2.2pt);
\foreach \v in {-1,1,2,3,4,5}{\node[below=4pt] at (\v,-1.5){$\v$};}\end{scope}
\node at (7.1,6.85) {$f(x)=\sin(\pi x),\quad a=1$};
\node[align=center] at (7.1,1.3) {$f(x+2)=f(x)$；$2$ 是一个正周期。};
\node[text width=12.6cm,align=center] at (7.1,.65) {橙点的横坐标为 $-\frac14,\frac14,\frac94,\frac{17}4$；曲线仅是满足全部条件的一个例子。};
\end{tikzpicture}\end{center}

奇函数本来就关于原点对称。点 $(x,y)$ 先绕原点转到 $(-x,-y)$，再绕 $(a,0)$ 转到 $(2a+x,y)$。高度先变成 $-y$，再变回 $y$，所以移动 $2|a|$ 已经重复。
\end{explanationbox}
